\subsection*{全体概要：この章で学ぶこと}

\addcontentsline{toc}{subsection}{全体概要：この章で学ぶこと}

この章は、ソフトウェア工学を専門職として実践するために必要な知識、技能、態度を扱う。ここでいう専門職実務とは、単にプログラムを書く能力ではない。倫理、標準、法的論点、契約、文書化、チームでの働き方、利害関係者との関係、読み書きや発表の技能を含む総合的な実務能力である。

ソフトウェアは、金融、医療、交通、教育、行政、娯楽、家庭内機器など、社会生活の広い範囲に入り込んでいる。したがって、ソフトウェア技術者の判断は、個人の利便性だけでなく、安全、プライバシー、財産、組織活動、社会的信頼にも影響する。第14章は、このような影響を理解し、責任ある態度で仕事を進めるための土台を示す章である。

\par\smallskip
\begin{center}
\centering
\IfFileExists{assets/generated_figures/ch14/fig-ch14-01.png}{\includegraphics[width=0.96\linewidth]{assets/generated_figures/ch14/fig-ch14-01.png}}{\fbox{\parbox[c][48mm][c]{0.9\linewidth}{\centering 画像生成待ち\\14章の全体地図}}}
\par\smallskip
\refstepcounter{figure}\label{fig:ch14-01}図 \thefigure: 14章の全体地図
\end{center}
図 \thefigure は、14章の全体地図を視覚的に整理したものである。
\par\smallskip
